\documentclass{beamer}
\usetheme{Madrid}
\usecolortheme{default}
\usepackage[utf8]{inputenc}
\usepackage{amsmath}
\usepackage{amsfonts}
\usepackage{amssymb}
\usepackage{tikz}

\title{CRISC Domain C - 100 More Very Hard Questions}
\subtitle{Risk Response and Reporting - Expert Level (Part 2)}
\author{CRISC Mastery Series}
\date{\today}

\begin{document}
	
	\begin{frame}
		\titlepage
	\end{frame}
	
	\begin{frame}{Outline}
		\tableofcontents
	\end{frame}
	
	% ============================================================
	% SECTION 1: RISK RESPONSE STRATEGIES - ADVANCED
	% ============================================================
	
	\section{1. Risk Response Strategies - Advanced}
	
	\begin{frame}{Q101: Risk Response Combination Trap}
		\begin{block}{Question}
			A company faces a \$3M risk from ransomware. They implement:
			\begin{itemize}
				\item Backups costing \$200K (reduces impact to \$500K)
				\item Cyber insurance costing \$150K (covers \$1M)
				\item Employee training costing \$50K (reduces likelihood by 40\%)
			\end{itemize}
			What is the NET risk exposure after all controls?
		\end{block}
		\begin{enumerate}[(A)]
			\item \$0
			\item \$500K
			\item \$250K (after insurance covers \$250K of remaining)
			\item \$1.5M
		\end{enumerate}
	\end{frame}
	
	\begin{frame}{Q101: Explanation}
		\begin{block}{Correct Answer: C}
			\textbf{Calculation:}
			\begin{itemize}
				\item Original risk: \$3M
				\item After backups: \$500K impact (but likelihood reduced by training)
				\item Training reduces likelihood 40\%: \$500K × 60\% = \$300K expected loss
				\item Insurance covers \$1M: covers remaining \$300K
				\item Net exposure: \$0 financial, but reputational risk remains
			\end{itemize}
		\end{block}
		\begin{alertblock}{Critical Trap}
			Multiple controls compound. Calculate sequentially, not additively.
		\end{alertblock}
	\end{frame}
	
	\begin{frame}{Q102: Risk Avoidance Hidden Cost Trap}
		\begin{block}{Question}
			A company avoids a risk by not entering a new market. What hidden risk does this create?
		\end{block}
		\begin{enumerate}[(A)]
			\item No hidden risk - avoidance is complete
			\item Opportunity cost and competitive disadvantage
			\item Increased regulatory scrutiny
			\item Higher insurance premiums
		\end{enumerate}
	\end{frame}
	
	\begin{frame}{Q102: Explanation}
		\begin{block}{Correct Answer: B}
			\textbf{Avoidance creates opportunity cost risk:}
			\begin{itemize}
				\item Lost revenue from not entering market
				\item Competitors may gain advantage
				\item Shareholders may question strategy
				\item Avoidance is not "free" - it has costs
			\end{itemize}
		\end{block}
		\begin{alertblock}{Critical Trap}
			Avoidance = Eliminates one risk but creates opportunity cost risk.
		\end{alertblock}
	\end{frame}
	
	\begin{frame}{Q103: Transfer of Unknown Risk}
		\begin{block}{Question}
			An organization purchases cyber insurance without conducting a risk assessment. What is the PRIMARY problem?
		\end{block}
		\begin{enumerate}[(A)]
			\item Insurance is too expensive
			\item Coverage may be inadequate or inappropriate
			\item The insurance company may deny claims
			\item Both B and C
		\end{enumerate}
	\end{frame}
	
	\begin{frame}{Q103: Explanation}
		\begin{block}{Correct Answer: D}
			\textbf{Without risk assessment:}
			\begin{itemize}
				\item You don't know what to insure
				\item Coverage limits may be wrong
				\item Exclusions may apply to your actual risks
				\item Insurer may deny claims for undisclosed risks
				\item Premiums may not reflect actual exposure
			\end{itemize}
		\end{block}
		\begin{alertblock}{Critical Trap}
			Transfer without assessment = Blind transfer = Potential coverage gap.
		\end{alertblock}
	\end{frame}
	
	\begin{frame}{Q104: Risk Acceptance Authority Escalation}
		\begin{block}{Question}
			A business unit manager accepts a risk with \$500K impact. The enterprise risk appetite is \$250K. What should the risk manager do?
		\end{block}
		\begin{enumerate}[(A)]
			\item Accept the decision - BU manager has authority
			\item Document the acceptance and move on
			\item Escalate to senior management for approval
			\item Automatically implement controls
		\end{enumerate}
	\end{frame}
	
	\begin{frame}{Q104: Explanation}
		\begin{block}{Correct Answer: C}
			\textbf{Escalation is required:}
			\begin{itemize}
				\item Risk exceeds enterprise appetite (\$500K > \$250K)
				\item BU manager cannot unilaterally exceed appetite
				\item Senior management must approve
				\item Board may need to be informed
			\end{itemize}
		\end{block}
		\begin{alertblock}{Critical Trap}
			Risk appetite is an enterprise limit, not a guideline.
			Exceeding it requires executive approval.
		\end{alertblock}
	\end{frame}
	
	\begin{frame}{Q105: Mitigation Timing Trap}
		\begin{block}{Question}
			A risk has 90\% likelihood and \$1M impact. Controls cost \$100K and take 6 months to implement. What is the risk of waiting 6 months?
		\end{block}
		\begin{enumerate}[(A)]
			\item No risk - controls will be implemented
			\item \$450K expected loss during implementation period
			\item \$900K expected loss during implementation period
			\item The risk will decrease over time
		\end{enumerate}
	\end{frame}
	
	\begin{frame}{Q105: Explanation}
		\begin{block}{Correct Answer: B}
			\textbf{Calculation:}
			\begin{itemize}
				\item Expected loss = 90\% × \$1M = \$900K annually
				\item During 6 months: \$900K × 0.5 = \$450K
				\item Risk continues while controls are being implemented
			\end{itemize}
		\end{block}
		\begin{alertblock}{Critical Trap}
			Implementation time = Risk exposure continues.
			Consider interim controls during implementation.
		\end{alertblock}
	\end{frame}
	
	\begin{frame}{Q106: Residual Risk Acceptance}
		\begin{block}{Question}
			After implementing controls, residual risk is \$400K. Risk appetite is \$250K. What are the options?
		\end{block}
		\begin{enumerate}[(A)]
			\item Accept residual risk - controls are implemented
			\item Implement additional controls, transfer, or escalate for acceptance
			\item Ignore the gap - close to appetite
			\item Remove all controls and start over
		\end{enumerate}
	\end{frame}
	
	\begin{frame}{Q106: Explanation}
		\begin{block}{Correct Answer: B}
			\textbf{Residual risk exceeds appetite:}
			\begin{itemize}
				\item Additional controls to reduce further
				\item Transfer remaining risk (insurance)
				\item Escalate to management for formal acceptance
				\item Cannot simply accept without authority
			\end{itemize}
		\end{block}
		\begin{alertblock}{Critical Trap}
			Residual risk > Appetite = Action required.
			Cannot ignore the gap.
		\end{alertblock}
	\end{frame}
	
	\begin{frame}{Q107: Transfer Limitation Trap}
		\begin{block}{Question}
			A company transfers risk through outsourcing. The vendor suffers a breach. Who bears the reputational damage?
		\end{block}
		\begin{enumerate}[(A)]
			\item The vendor
			\item The original company
			\item Both equally
			\item No one - reputational damage is transferred
		\end{enumerate}
	\end{frame}
	
	\begin{frame}{Q107: Explanation}
		\begin{block}{Correct Answer: B}
			\textbf{Reputational damage remains with the original company:}
			\begin{itemize}
				\item Customers blame the brand they trusted
				\item Media reports name the original company
				\item Regulatory scrutiny focuses on the data owner
				\item Reputation cannot be outsourced
			\end{itemize}
		\end{block}
		\begin{alertblock}{Critical Trap}
			Transfer = Financial only.
			Reputation NEVER transfers.
		\end{alertblock}
	\end{frame}
	
	\begin{frame}{Q108: Avoidance Decision Trap}
		\begin{block}{Question}
			A company avoids a risk by terminating a product line. What must be considered before this decision?
		\end{block}
		\begin{enumerate}[(A)]
			\item Only the risk being avoided
			\item The risk, opportunity cost, and stakeholder impact
			\item Only the cost of termination
			\item Only regulatory requirements
		\end{enumerate}
	\end{frame}
	
	\begin{frame}{Q108: Explanation}
		\begin{block}{Correct Answer: B}
			\textbf{Comprehensive analysis required:}
			\begin{itemize}
				\item What risk is being avoided?
				\item What revenue is lost?
				\item How are customers affected?
				\item What are employee impacts?
				\item What is the competitive impact?
			\end{itemize}
		\end{block}
		\begin{alertblock}{Critical Trap}
			Avoidance is irreversible - consider all impacts before deciding.
		\end{alertblock}
	\end{frame}
	
	\begin{frame}{Q109: Risk Response Timing Trap}
		\begin{block}{Question}
			A risk is identified with 50\% likelihood and \$100K impact. The control costs \$60K. Should the control be implemented?
		\end{block}
		\begin{enumerate}[(A)]
			\item Yes, because risk exists
			\item No, because control cost exceeds expected loss
			\item Yes, because 50\% is high likelihood
			\item Cannot be determined
		\end{enumerate}
	\end{frame}
	
	\begin{frame}{Q109: Explanation}
		\begin{block}{Correct Answer: B}
			\textbf{Calculation:}
			\begin{itemize}
				\item Expected loss = 50\% × \$100K = \$50K
				\item Control cost = \$60K
				\item \$60K > \$50K = Control NOT justified
				\item Accept the risk (documented)
			\end{itemize}
		\end{block}
		\begin{alertblock}{Critical Trap}
			Expected loss = Likelihood × Impact.
			Control cost > Expected loss = Accept.
		\end{alertblock}
	\end{frame}
	
	\begin{frame}{Q110: Risk Response Documentation Trap}
		\begin{block}{Question}
			A risk response was implemented but not documented. What is the consequence?
		\end{block}
		\begin{enumerate}[(A)]
			\item No consequence - the risk is mitigated
			\item Cannot demonstrate due care or effectiveness
			\item The control automatically fails
			\item The risk increases
		\end{enumerate}
	\end{frame}
	
	\begin{frame}{Q110: Explanation}
		\begin{block}{Correct Answer: B}
			\textbf{Undocumented controls have limited value:}
			\begin{itemize}
				\item Cannot prove control exists
				\item Cannot demonstrate effectiveness
				\item Cannot defend in audit or legal proceedings
				\item Knowledge is lost when personnel change
			\end{itemize}
		\end{block}
		\begin{alertblock}{Critical Trap}
			Undocumented = Undefensible.
			Document all risk responses.
		\end{alertblock}
	\end{frame}
	
	% ============================================================
	% SECTION 2: RISK AND CONTROL OWNERSHIP - ADVANCED
	% ============================================================
	
	\section{2. Risk and Control Ownership - Advanced}
	
	\begin{frame}{Q111: Risk Owner Conflict Resolution}
		\begin{block}{Question}
			Two business units claim ownership of a shared risk. Neither wants to take responsibility. What should happen?
		\end{block}
		\begin{enumerate}[(A)]
			\item Split the risk 50/50
			\item Assign to the BU with larger budget
			\item Escalate to senior management to assign a single owner
			\item The risk remains unowned
		\end{enumerate}
	\end{frame}
	
	\begin{frame}{Q111: Explanation}
		\begin{block}{Correct Answer: C}
			\textbf{Escalate for single ownership:}
			\begin{itemize}
				\item Shared ownership = No accountability
				\item Senior management can assign authority
				\item Owner must have cross-functional authority
				\item May need C-level or COO assignment
			\end{itemize}
		\end{block}
		\begin{alertblock}{Critical Trap}
			Multiple owners = No owner.
			Single point of accountability required.
		\end{alertblock}
	\end{frame}
	
	\begin{frame}{Q112: Control Owner Without Authority Trap}
		\begin{block}{Question}
			A control owner identifies a failing control but has no budget to fix it. What should they do FIRST?
		\end{block}
		\begin{enumerate}[(A)]
			\item Document the issue and escalate to risk owner
			\item Implement a compensating control themselves
			\item Accept the risk and move on
			\item Wait for the next audit to report
		\end{enumerate}
	\end{frame}
	
	\begin{frame}{Q112: Explanation}
		\begin{block}{Correct Answer: A}
			\textbf{Escalate to risk owner:}
			\begin{itemize}
				\item Risk owner has budget authority
				\item Risk owner is accountable
				\item Control owner documents the issue
				\item Control owner may recommend solutions
			\end{itemize}
		\end{block}
		\begin{alertblock}{Critical Trap}
			Control owner not equal to Budget authority.
			Escalate to risk owner for resources.
		\end{alertblock}
	\end{frame}
	
	\begin{frame}{Q113: Risk Owner Delegation Trap}
		\begin{block}{Question}
			A risk owner delegates authority to a subordinate without documentation. What is the risk?
		\end{block}
		\begin{enumerate}[(A)]
			\item No risk - delegation is informal
			\item Subordinate cannot be held accountable
			\item Risk owner remains accountable but subordinate has no authority
			\item Both B and C
		\end{enumerate}
	\end{frame}
	
	\begin{frame}{Q113: Explanation}
		\begin{block}{Correct Answer: D}
			\textbf{Undocumented delegation creates issues:}
			\begin{itemize}
				\item Subordinate cannot make binding decisions
				\item Risk owner remains accountable
				\item No clear authority chain
				\item Audit trail is broken
			\end{itemize}
		\end{block}
		\begin{alertblock}{Critical Trap}
			Delegation must be documented, communicated, and bounded.
		\end{alertblock}
	\end{frame}
	
	\begin{frame}{Q114: Control Owner Multiple Controls}
		\begin{block}{Question}
			A control owner is responsible for 50 controls. What is the PRIMARY risk?
		\end{block}
		\begin{enumerate}[(A)]
			\item Too much work
			\item Controls may not receive adequate attention
			\item Higher salary required
			\item Needs more training
		\end{enumerate}
	\end{frame}
	
	\begin{frame}{Q114: Explanation}
		\begin{block}{Correct Answer: B}
			\textbf{Span of control risk:}
			\begin{itemize}
				\item 50 controls is too many for one person
				\item Some controls will be neglected
				\item Testing and monitoring may be insufficient
				\item Effectiveness may degrade
			\end{itemize}
		\end{block}
		\begin{alertblock}{Critical Trap}
			Too many controls = Diluted oversight.
			Consider control owner workload.
		\end{alertblock}
	\end{frame}
	
	\begin{frame}{Q115: Risk Owner Out of Office}
		\begin{block}{Question}
			A critical risk decision is needed but the risk owner is on vacation. What should happen?
		\end{block}
		\begin{enumerate}[(A)]
			\item Wait for the risk owner to return
			\item The control owner can make the decision
			\item A designated alternate with documented authority
			\item The decision cannot be made
		\end{enumerate}
	\end{frame}
	
	\begin{frame}{Q115: Explanation}
		\begin{block}{Correct Answer: C}
			\textbf{Designated alternate required:}
			\begin{itemize}
				\item Risk owner should designate an alternate
				\item Alternate must have documented authority
				\item Alternate must understand the risk
				\item Business continuity for risk management
			\end{itemize}
		\end{block}
		\begin{alertblock}{Critical Trap}
			No alternate = Decision paralysis.
			Always designate an alternate risk owner.
		\end{alertblock}
	\end{frame}
	
	\begin{frame}{Q116: Risk Owner vs Control Owner Conflict}
		\begin{block}{Question}
			Risk owner wants to accept risk. Control owner insists on remediation. Who decides?
		\end{block}
		\begin{enumerate}[(A)]
			\item Risk owner - they own the risk
			\item Control owner - they own the control
			\item Senior management - conflict requires escalation
			\item Whoever has more budget
		\end{enumerate}
	\end{frame}
	
	\begin{frame}{Q116: Explanation}
		\begin{block}{Correct Answer: C}
			\textbf{Escalate to senior management:}
			\begin{itemize}
				\item Conflict indicates unresolved issues
				\item Senior management provides perspective
				\item Balance risk and control needs
				\item Document the decision
			\end{itemize}
		\end{block}
		\begin{alertblock}{Critical Trap}
			Risk owner not equal to Unilateral authority.
			Conflicts require escalation.
		\end{alertblock}
	\end{frame}
	
	\begin{frame}{Q117: Risk Owner for IT Risk}
		\begin{block}{Question}
			Who should be the risk owner for a data breach risk?
		\end{block}
		\begin{enumerate}[(A)]
			\item CISO
			\item CIO
			\item Business leader who owns the data
			\item IT manager
		\end{enumerate}
	\end{frame}
	
	\begin{frame}{Q117: Explanation}
		\begin{block}{Correct Answer: C}
			\textbf{Business leader owns the risk:}
			\begin{itemize}
				\item Data is a business asset
				\item Breach impacts business objectives
				\item Business leader has budget authority
				\item IT/CISO are enablers, not owners
			\end{itemize}
		\end{block}
		\begin{alertblock}{Critical Trap}
			IT does NOT own business risk.
			Data owner = Risk owner.
		\end{alertblock}
	\end{frame}
	
	\begin{frame}{Q118: Control Owner Testing Responsibility}
		\begin{block}{Question}
			Who is responsible for testing controls?
		\end{block}
		\begin{enumerate}[(A)]
			\item Risk owner
			\item Control owner
			\item Internal audit
			\item External auditor
		\end{enumerate}
	\end{frame}
	
	\begin{frame}{Q118: Explanation}
		\begin{block}{Correct Answer: B}
			\textbf{Control owner tests controls:}
			\begin{itemize}
				\item Control owner is responsible for effectiveness
				\item Testing verifies the control works
				\item Internal audit provides independent assurance
				\item External audit provides external assurance
			\end{itemize}
		\end{block}
		\begin{alertblock}{Critical Trap}
			Control owner = First line of testing.
			Audit = Independent verification.
		\end{alertblock}
	\end{frame}
	
	\begin{frame}{Q119: Risk Owner Budget Authority}
		\begin{block}{Question}
			A risk owner lacks budget to implement controls. What is the consequence?
		\end{block}
		\begin{enumerate}[(A)]
			\item Nothing - controls are optional
			\item Risk cannot be effectively managed
			\item Control owner must pay
			\item Risk is automatically transferred
		\end{enumerate}
	\end{frame}
	
	\begin{frame}{Q119: Explanation}
		\begin{block}{Correct Answer: B}
			\textbf{No budget = No effective risk management:}
			\begin{itemize}
				\item Controls cannot be implemented
				\item Risk remains unmanaged
				\item Risk owner is ineffective
				\item Escalation to higher authority needed
			\end{itemize}
		\end{block}
		\begin{alertblock}{Critical Trap}
			Risk ownership requires budget authority.
			No budget = Not a true risk owner.
		\end{alertblock}
	\end{frame}
	
	\begin{frame}{Q120: Risk Owner Rotation}
		\begin{block}{Question}
			Should risk ownership rotate periodically?
		\end{block}
		\begin{enumerate}[(A)]
			\item Yes, for objectivity
			\item No, ownership should be stable
			\item Depends on organizational requirements
			\item Only for IT risks
		\end{enumerate}
	\end{frame}
	
	\begin{frame}{Q120: Explanation}
		\begin{block}{Correct Answer: B}
			\textbf{Risk ownership should be stable:}
			\begin{itemize}
				\item Continuity of knowledge
				\item Accountability is maintained
				\item Relationships with stakeholders
				\item Long-term risk perspective
			\end{itemize}
		\end{block}
		\begin{alertblock}{Critical Trap}
			Risk owner turnover = Loss of continuity.
			Stability is preferred.
		\end{alertblock}
	\end{frame}
	
	% ============================================================
	% SECTION 3: THIRD-PARTY RISK MANAGEMENT - ADVANCED
	% ============================================================
	
	\section{3. Third-Party Risk Management - Advanced}
	
	\begin{frame}{Q121: Vendor Assessment Scope Trap}
		\begin{block}{Question}
			A vendor has a SOC 2 Type II report. Can you rely on it for all services?
		\end{block}
		\begin{enumerate}[(A)]
			\item Yes, SOC 2 covers all services
			\item No, only services within the report scope
			\item Yes, if the report is recent
			\item No, SOC 2 is not reliable
		\end{enumerate}
	\end{frame}
	
	\begin{frame}{Q121: Explanation}
		\begin{block}{Correct Answer: B}
			\textbf{SOC 2 scope is limited:}
			\begin{itemize}
				\item Report covers specific services
				\item Out-of-scope services need assessment
				\item Check the scope section first
				\item Request scope expansion if needed
			\end{itemize}
		\end{block}
		\begin{alertblock}{Critical Trap}
			SOC 2 not equal to All services.
			Scope matters.
		\end{alertblock}
	\end{frame}
	
	\begin{frame}{Q122: Fourth-Party Risk Trap}
		\begin{block}{Question}
			Your vendor uses a subcontractor. The subcontractor has a breach. Who is liable?
		\end{block}
		\begin{enumerate}[(A)]
			\item Only the subcontractor
			\item Your organization and the vendor
			\item Only your organization
			\item Only the vendor
		\end{enumerate}
	\end{frame}
	
	\begin{frame}{Q122: Explanation}
		\begin{block}{Correct Answer: B}
			\textbf{Multiple parties liable:}
			\begin{itemize}
				\item Your organization: Regulatory accountability
				\item Vendor: Contractual responsibility
				\item Subcontractor: Direct liability
				\item Fourth-party risk affects you
			\end{itemize}
		\end{block}
		\begin{alertblock}{Critical Trap}
			Fourth-party breach = Your problem too.
			You cannot outsource accountability.
		\end{alertblock}
	\end{frame}
	
	\begin{frame}{Q123: Vendor Offboarding Trap}
		\begin{block}{Question}
			What is the GREATEST risk during vendor offboarding?
		\end{block}
		\begin{enumerate}[(A)]
			\item Service disruption
			\item Data remanence (data not destroyed)
			\item Cost of transition
			\item Loss of knowledge
		\end{enumerate}
	\end{frame}
	
	\begin{frame}{Q123: Explanation}
		\begin{block}{Correct Answer: B}
			\textbf{Data remanence is greatest risk:}
			\begin{itemize}
				\item Data may remain on vendor systems
				\item Backups may not be deleted
				\item Regulatory violations possible
				\item Reputational damage if leaked
			\end{itemize}
		\end{block}
		\begin{alertblock}{Critical Trap}
			Offboarding = Data destruction is 1 priority.
			Require certificate of destruction.
		\end{alertblock}
	\end{frame}
	
	\begin{frame}{Q124: Vendor Contract Trap}
		\begin{block}{Question}
			Which contract clause is MOST important for data protection?
		\end{block}
		\begin{enumerate}[(A)]
			\item Service level agreement
			\item Right to audit and security requirements
			\item Indemnification clause
			\item Non-disclosure agreement
		\end{enumerate}
	\end{frame}
	
	\begin{frame}{Q124: Explanation}
		\begin{block}{Correct Answer: B}
			\textbf{Right to audit + security requirements:}
			\begin{itemize}
				\item Enables ongoing verification
				\item Establishes minimum standards
				\item Provides recourse if standards not met
				\item Financial penalties don't prevent breaches
			\end{itemize}
		\end{block}
		\begin{alertblock}{Critical Trap}
			Prevention > Compensation.
			Right to audit enables prevention.
		\end{alertblock}
	\end{frame}
	
	\begin{frame}{Q125: Vendor Assessment Timing Trap}
		\begin{block}{Question}
			When should a vendor assessment be conducted?
		\end{block}
		\begin{enumerate}[(A)]
			\item Before contract signing
			\item After contract signing
			\item After go-live
			\item Annually only
		\end{enumerate}
	\end{frame}
	
	\begin{frame}{Q125: Explanation}
		\begin{block}{Correct Answer: A}
			\textbf{Before contract signing:}
			\begin{itemize}
				\item Leverage to negotiate security terms
				\item Informed decision about vendor
				\item Identify deal-breakers early
				\item Can require improvements before signing
			\end{itemize}
		\end{block}
		\begin{alertblock}{Critical Trap}
			Assessment after contract = No leverage.
			Assess before signing.
		\end{alertblock}
	\end{frame}
	
	\begin{frame}{Q126: Vendor Monitoring Trap}
		\begin{block}{Question}
			How often should vendor risk be reassessed?
		\end{block}
		\begin{enumerate}[(A)]
			\item Only at contract renewal
			\item Annually or when significant changes occur
			\item Every 3 years
			\item Only when there is a breach
		\end{enumerate}
	\end{frame}
	
	\begin{frame}{Q126: Explanation}
		\begin{block}{Correct Answer: B}
			\textbf{Annual or event-driven reassessment:}
			\begin{itemize}
				\item Annual review minimum
				\item When vendor services change
				\item When vendor has incident
				\item When regulations change
				\item When vendor's security posture changes
			\end{itemize}
		\end{block}
		\begin{alertblock}{Critical Trap}
			Vendor risk is not static.
			Annual reassessment required.
		\end{alertblock}
	\end{frame}
	
	\begin{frame}{Q127: Vendor Exception Management}
		\begin{block}{Question}
			A vendor requests an exception to your security policy. What should happen?
		\end{block}
		\begin{enumerate}[(A)]
			\item Approve automatically - vendor knows best
			\item Evaluate risk and require compensating controls
			\item Deny all exceptions
			\item Ignore the request
		\end{enumerate}
	\end{frame}
	
	\begin{frame}{Q127: Explanation}
		\begin{block}{Correct Answer: B}
			\textbf{Evaluate risk and require compensating controls:}
			\begin{itemize}
				\item Understand the business need
				\item Assess the risk of the exception
				\item Require compensating controls
				\item Document and review periodically
			\end{itemize}
		\end{block}
		\begin{alertblock}{Critical Trap}
			Exceptions are not automatic.
			Evaluate, compensate, document.
		\end{alertblock}
	\end{frame}
	
	\begin{frame}{Q128: Vendor Incident Notification}
		\begin{block}{Question}
			How quickly should a vendor notify you of a breach?
		\end{block}
		\begin{enumerate}[(A)]
			\item Within 30 days
			\item Within 24-72 hours (per contract/regulation)
			\item Within 1 week
			\item At their discretion
		\end{enumerate}
	\end{frame}
	
	\begin{frame}{Q128: Explanation}
		\begin{block}{Correct Answer: B}
			\textbf{24-72 hours is standard:}
			\begin{itemize}
				\item Regulatory requirements (GDPR: 72 hours)
				\item Contract should specify timeframe
				\item You need time to respond
				\item You need to notify regulators
			\end{itemize}
		\end{block}
		\begin{alertblock}{Critical Trap}
			Delayed notification = Delayed response.
			Specify notification timeframe in contract.
		\end{alertblock}
	\end{frame}
	
	\begin{frame}{Q129: Vendor Data Location Trap}
		\begin{block}{Question}
			A vendor stores your data in a foreign country. What is the PRIMARY concern?
		\end{block}
		\begin{enumerate}[(A)]
			\item Data latency
			\item Regulatory and legal jurisdiction issues
			\item Higher costs
			\item Language barriers
		\end{enumerate}
	\end{frame}
	
	\begin{frame}{Q129: Explanation}
		\begin{block}{Correct Answer: B}
			\textbf{Jurisdictional issues:}
			\begin{itemize}
				\item Data protection laws vary by country
				\item GDPR, CCPA may apply
				\item Law enforcement access differs
				\item Legal process may be difficult
			\end{itemize}
		\end{block}
		\begin{alertblock}{Critical Trap}
			Data location = Legal jurisdiction.
			Know where your data is stored.
		\end{alertblock}
	\end{frame}
	
	\begin{frame}{Q130: Vendor Subcontractor Approval}
		\begin{block}{Question}
			Should a vendor be required to obtain approval before using subcontractors?
		\end{block}
		\begin{enumerate}[(A)]
			\item No, vendors can manage their own operations
			\item Yes, to maintain visibility of fourth-party risk
			\item Only if the subcontractor is in a different country
			\item Only for critical services
		\end{enumerate}
	\end{frame}
	
	\begin{frame}{Q130: Explanation}
		\begin{block}{Correct Answer: B}
			\textbf{Approval required for visibility:}
			\begin{itemize}
				\item You need to know who handles your data
				\item Fourth-party risk must be assessed
				\item You may object to specific subcontractors
				\item Maintains chain of accountability
			\end{itemize}
		\end{block}
		\begin{alertblock}{Critical Trap}
			Unknown subcontractors = Unmanaged risk.
			Require approval for subcontractors.
		\end{alertblock}
	\end{frame}
	
	% ============================================================
	% SECTION 4: CONTROL DESIGN AND IMPLEMENTATION - ADVANCED
	% ============================================================
	
	\section{4. Control Design and Implementation - Advanced}
	
	\begin{frame}{Q131: Preventive vs Deterrent Trap}
		\begin{block}{Question}
			A warning banner on a login page is what type of control?
		\end{block}
		\begin{enumerate}[(A)]
			\item Preventive
			\item Detective
			\item Deterrent
			\item Corrective
		\end{enumerate}
	\end{frame}
	
	\begin{frame}{Q131: Explanation}
		\begin{block}{Correct Answer: C}
			\textbf{Warning banner is deterrent:}
			\begin{itemize}
				\item It warns of consequences
				\item It doesn't block access (preventive)
				\item It doesn't detect intrusion (detective)
				\item It discourages unauthorized access through fear
			\end{itemize}
		\end{block}
		\begin{alertblock}{Critical Trap}
			Warning = Deterrent.
			It discourages, it doesn't prevent.
		\end{alertblock}
	\end{frame}
	
	\begin{frame}{Q132: Detective vs Corrective Trap}
		\begin{block}{Question}
			An IDS detects and blocks an intrusion. Is this detective or corrective?
		\end{block}
		\begin{enumerate}[(A)]
			\item Detective only
			\item Corrective only
			\item Both detective and corrective
			\item Preventive
		\end{enumerate}
	\end{frame}
	
	\begin{frame}{Q132: Explanation}
		\begin{block}{Correct Answer: C}
			\textbf{IDS with auto-block performs both:}
			\begin{itemize}
				\item Detection = Detective
				\item Blocking = Corrective
				\item Modern systems combine functions
				\item Analyze the actions, not the name
			\end{itemize}
		\end{block}
		\begin{alertblock}{Critical Trap}
			Controls can have multiple functions.
			IDS with auto-block = Detective + Corrective.
		\end{alertblock}
	\end{frame}
	
	\begin{frame}{Q133: Compensating Control Trap}
		\begin{block}{Question}
			A compensating control must provide:
		\end{block}
		\begin{enumerate}[(A)]
			\item Cheaper solution
			\item Equivalent or superior risk reduction
			\item Automated solution
			\item IT-approved solution
		\end{enumerate}
	\end{frame}
	
	\begin{frame}{Q133: Explanation}
		\begin{block}{Correct Answer: B}
			\textbf{Equivalent or superior risk reduction:}
			\begin{itemize}
				\item Compensating control must address the same risk
				\item Must reduce risk to acceptable level
				\item Cost is secondary
				\item Effectiveness is primary
			\end{itemize}
		\end{block}
		\begin{alertblock}{Critical Trap}
			Compensating control = Same risk reduction, different method.
			Cheaper but less effective = Not acceptable.
		\end{alertblock}
	\end{frame}
	
	\begin{frame}{Q134: Control Testing Frequency}
		\begin{block}{Question}
			How often should controls be tested?
		\end{block}
		\begin{enumerate}[(A)]
			\item Annually for all controls
			\item Based on risk and criticality
			\item Only when changes occur
			\item Quarterly for all controls
		\end{enumerate}
	\end{frame}
	
	\begin{frame}{Q134: Explanation}
		\begin{block}{Correct Answer: B}
			\textbf{Risk-based testing frequency:}
			\begin{itemize}
				\item High-risk controls: More frequent
				\item Low-risk controls: Less frequent
				\item Regulatory requirements may mandate
				\item Consider change frequency and stability
			\end{itemize}
		\end{block}
		\begin{alertblock}{Critical Trap}
			One size does NOT fit all.
			Risk-based frequency is appropriate.
		\end{alertblock}
	\end{frame}
	
	\begin{frame}{Q135: Design vs Operating Deficiency}
		\begin{block}{Question}
			A control is designed correctly but not operating as designed. This is:
		\end{block}
		\begin{enumerate}[(A)]
			\item Design deficiency
			\item Operating deficiency
			\item Control gap
			\item Control redundancy
		\end{enumerate}
	\end{frame}
	
	\begin{frame}{Q135: Explanation}
		\begin{block}{Correct Answer: B}
			\textbf{Operating deficiency:}
			\begin{itemize}
				\item Control is designed correctly
				\item Control is not operating effectively
				\item Could be due to personnel, process, or technology failure
				\item Remediation focuses on operation
			\end{itemize}
		\end{block}
		\begin{alertblock}{Critical Trap}
			Design = Blueprint.
			Operating = Actual performance.
			Design correct, operation failing = Operating deficiency.
		\end{alertblock}
	\end{frame}
	
	\begin{frame}{Q136: Parallel Changeover Risk}
		\begin{block}{Question}
			What is the PRIMARY risk of parallel changeover?
		\end{block}
		\begin{enumerate}[(A)]
			\item System failure
			\item Data inconsistency
			\item User resistance
			\item Higher costs
		\end{enumerate}
	\end{frame}
	
	\begin{frame}{Q136: Explanation}
		\begin{block}{Correct Answer: B}
			\textbf{Data inconsistency is primary risk:}
			\begin{itemize}
				\item Both systems process transactions
				\item Data may not sync
				\item Reconciliation is complex
				\item Errors may go undetected
			\end{itemize}
		\end{block}
		\begin{alertblock}{Critical Trap}
			Parallel = Safest for availability, riskiest for data consistency.
		\end{alertblock}
	\end{frame}
	
	\begin{frame}{Q137: Abrupt Changeover Appropriate}
		\begin{block}{Question}
			When is abrupt changeover acceptable?
		\end{block}
		\begin{enumerate}[(A)]
			\item For critical systems
			\item For minor changes with easy rollback
			\item When users are trained
			\item When vendor recommends
		\end{enumerate}
	\end{frame}
	
	\begin{frame}{Q137: Explanation}
		\begin{block}{Correct Answer: B}
			\textbf{Minor changes with easy rollback:}
			\begin{itemize}
				\item Low business impact
				\item Quick recovery possible
				\item No critical processes affected
				\item Acceptable risk level
			\end{itemize}
		\end{block}
		\begin{alertblock}{Critical Trap}
			Abrupt = Fastest but highest risk.
			Only for non-critical, reversible changes.
		\end{alertblock}
	\end{frame}
	
	\begin{frame}{Q138: Post-Implementation Review Timing}
		\begin{block}{Question}
			When should a post-implementation review occur?
		\end{block}
		\begin{enumerate}[(A)]
			\item Immediately after go-live
			\item After system has stabilized
			\item Only if issues occurred
			\item Before next audit
		\end{enumerate}
	\end{frame}
	
	\begin{frame}{Q138: Explanation}
		\begin{block}{Correct Answer: B}
			\textbf{After stabilization (3-6 months):}
			\begin{itemize}
				\item Issues have surfaced
				\item User feedback is meaningful
				\item Performance can be assessed
				\item Lessons learned are clear
			\end{itemize}
		\end{block}
		\begin{alertblock}{Critical Trap}
			Too early = Issues haven't surfaced.
			Too late = Lessons forgotten.
			Wait for stability.
		\end{alertblock}
	\end{frame}
	
	\begin{frame}{Q139: Control Automation Trap}
		\begin{block}{Question}
			A control is automated. Does this eliminate the need for monitoring?
		\end{block}
		\begin{enumerate}[(A)]
			\item Yes, automated controls are self-monitoring
			\item No, automated controls can fail
			\item Yes, automation eliminates errors
			\item No, but monitoring can be reduced
		\end{enumerate}
	\end{frame}
	
	\begin{frame}{Q139: Explanation}
		\begin{block}{Correct Answer: B}
			\textbf{Automated controls still require monitoring:}
			\begin{itemize}
				\item Automation can fail (bugs, config errors)
				\item Automated controls can be bypassed
				\item Automation logic may become outdated
				\item Regular testing confirms effectiveness
			\end{itemize}
		\end{block}
		\begin{alertblock}{Critical Trap}
			Automated not equal to Self-monitoring.
			Automation reduces human error but doesn't eliminate oversight.
		\end{alertblock}
	\end{frame}
	
	\begin{frame}{Q140: Control Testing Approach}
		\begin{block}{Question}
			What is the BEST approach for control testing?
		\end{block}
		\begin{enumerate}[(A)]
			\item Test every transaction
			\item Statistical sampling with documented methodology
			\item Test only recent transactions
			\item Rely on management assurance
		\end{enumerate}
	\end{frame}
	
	\begin{frame}{Q140: Explanation}
		\begin{block}{Correct Answer: B}
			\textbf{Statistical sampling is best:}
			\begin{itemize}
				\item Testing every transaction is impractical
				\item Sampling provides reasonable assurance
				\item Documented methodology enables repeatability
				\item Results are defensible to auditors
			\end{itemize}
		\end{block}
		\begin{alertblock}{Critical Trap}
			Testing = Sampling (not 100\%).
			Document methodology for auditability.
		\end{alertblock}
	\end{frame}
	
	% ============================================================
	% SECTION 5: MONITORING AND REPORTING - ADVANCED
	% ============================================================
	
	\section{5. Monitoring and Reporting - Advanced}
	
	\begin{frame}{Q141: KRI vs KPI vs KCI Trap}
		\begin{block}{Question}
			The percentage of employees who completed security awareness training is a:
		\end{block}
		\begin{enumerate}[(A)]
			\item KRI
			\item KPI
			\item KCI
			\item None
		\end{enumerate}
	\end{frame}
	
	\begin{frame}{Q141: Explanation}
		\begin{block}{Correct Answer: B}
			\textbf{Training completion is KPI:}
			\begin{itemize}
				\item Measures performance of training program
				\item Backward-looking
				\item Doesn't predict risk (KRI)
				\item Doesn't measure control effectiveness (KCI)
			\end{itemize}
		\end{block}
		\begin{alertblock}{Critical Trap}
			Training completion = KPI.
			Phishing click rate = KRI.
			Access review completion = KCI.
		\end{alertblock}
	\end{frame}
	
	\begin{frame}{Q142: KRI Threshold Breach}
		\begin{block}{Question}
			A KRI has breached its threshold. What should happen FIRST?
		\end{block}
		\begin{enumerate}[(A)]
			\item Escalate to senior management
			\item Investigate the cause
			\item Implement additional controls
			\item Update the risk register
		\end{enumerate}
	\end{frame}
	
	\begin{frame}{Q142: Explanation}
		\begin{block}{Correct Answer: B}
			\textbf{Investigate first:}
			\begin{itemize}
				\item Is the data accurate?
				\item Is it a spike or a trend?
				\item What is the root cause?
				\item What response is appropriate?
			\end{itemize}
		\end{block}
		\begin{alertblock}{Critical Trap}
			Threshold breach not equal to Immediate escalation.
			Investigate first, then respond.
		\end{alertblock}
	\end{frame}
	
	\begin{frame}{Q143: Board Reporting Trap}
		\begin{block}{Question}
			What should NOT be included in a board-level risk report?
		\end{block}
		\begin{enumerate}[(A)]
			\item Key risk indicators
			\item Comparison to risk appetite
			\item Technical vulnerability details
			\item Recommended actions
		\end{enumerate}
	\end{frame}
	
	\begin{frame}{Q143: Explanation}
		\begin{block}{Correct Answer: C}
			\textbf{Technical details should NOT be in board reports:}
			\begin{itemize}
				\item Board members are not technical
				\item Focus on business impact
				\item Technical details obscure message
				\item Board needs strategic view
			\end{itemize}
		\end{block}
		\begin{alertblock}{Critical Trap}
			Board = Business language.
			Technical details = Operational reports.
		\end{alertblock}
	\end{frame}
	
	\begin{frame}{Q144: SMART KRI Criteria}
		\begin{block}{Question}
			A KRI is defined as "number of security incidents." What is the weakness?
		\end{block}
		\begin{enumerate}[(A)]
			\item Not measurable
			\item Lacks specificity and thresholds
			\item Too difficult to track
			\item Not relevant
		\end{enumerate}
	\end{frame}
	
	\begin{frame}{Q144: Explanation}
		\begin{block}{Correct Answer: B}
			\textbf{Lacks specificity and thresholds:}
			\begin{itemize}
				\item What type of incidents?
				\item What severity?
				\item What timeframe?
				\item What triggers action?
			\end{itemize}
		\end{block}
		\begin{alertblock}{Critical Trap}
			Vague KRI = Useless KRI.
			Specific + Measurable + Threshold + Actionable.
		\end{alertblock}
	\end{frame}
	
	\begin{frame}{Q145: Control Monitoring Purpose}
		\begin{block}{Question}
			The PRIMARY purpose of control monitoring is to:
		\end{block}
		\begin{enumerate}[(A)]
			\item Verify control is working as designed
			\item Verify control is effectively reducing risk
			\item Verify control is cost-effective
			\item Verify control is automated
		\end{enumerate}
	\end{frame}
	
	\begin{frame}{Q145: Explanation}
		\begin{block}{Correct Answer: B}
			\textbf{Verify risk reduction effectiveness:}
			\begin{itemize}
				\item Control working not equal to Control effective
				\item Design may be flawed
				\item Threat landscape may have changed
				\item Monitor effectiveness, not just operation
			\end{itemize}
		\end{block}
		\begin{alertblock}{Critical Trap}
			Control working not equal to Control effective.
			Monitor effectiveness.
		\end{alertblock}
	\end{frame}
	
	\begin{frame}{Q146: Heat Map Limitation}
		\begin{block}{Question}
			What is the PRIMARY limitation of a risk heat map?
		\end{block}
		\begin{enumerate}[(A)]
			\item Too complex
			\item Qualitative and subjective
			\item Cannot show multiple risks
			\item Requires too much data
		\end{enumerate}
	\end{frame}
	
	\begin{frame}{Q146: Explanation}
		\begin{block}{Correct Answer: B}
			\textbf{Qualitative and subjective:}
			\begin{itemize}
				\item "High" means different things
				\item No quantification
				\item Difficult to compare
				\item Can oversimplify
			\end{itemize}
		\end{block}
		\begin{alertblock}{Critical Trap}
			Heat map = Qualitative = Subjective.
			Dashboards = Quantitative + Qualitative.
		\end{alertblock}
	\end{frame}
	
	\begin{frame}{Q147: Exception Report Content}
		\begin{block}{Question}
			What should an exception report include?
		\end{block}
		\begin{enumerate}[(A)]
			\item Only the exception list
			\item Exception, compensating controls, owner, expiration
			\item Only long-standing exceptions
			\item Only IT-approved exceptions
		\end{enumerate}
	\end{frame}
	
	\begin{frame}{Q147: Explanation}
		\begin{block}{Correct Answer: B}
			\textbf{Comprehensive exception reporting:}
			\begin{itemize}
				\item Description of exception
				\item Compensating controls
				\item Risk owner
				\item Expiration date
				\item Business justification
			\end{itemize}
		\end{block}
		\begin{alertblock}{Critical Trap}
			Exception report not equal to Simple list.
			Must include risk context.
		\end{alertblock}
	\end{frame}
	
	\begin{frame}{Q148: KCI vs KRI}
		\begin{block}{Question}
			The number of access reviews completed on time is a:
		\end{block}
		\begin{enumerate}[(A)]
			\item KRI
			\item KPI
			\item KCI
			\item None
		\end{enumerate}
	\end{frame}
	
	\begin{frame}{Q148: Explanation}
		\begin{block}{Correct Answer: C}
			\textbf{Access review completion is KCI:}
			\begin{itemize}
				\item Access review is a control
				\item Completion measures control effectiveness
				\item Doesn't predict risk (KRI)
				\item Doesn't measure overall performance (KPI)
			\end{itemize}
		\end{block}
		\begin{alertblock}{Critical Trap}
			If it measures a control = KCI.
			If it predicts risk = KRI.
			If it measures performance = KPI.
		\end{alertblock}
	\end{frame}
	
	\begin{frame}{Q149: Log Retention}
		\begin{block}{Question}
			How long should security logs be retained?
		\end{block}
		\begin{enumerate}[(A)]
			\item 30 days
			\item 90 days
			\item Based on regulatory and forensic needs
			\item Indefinitely
		\end{enumerate}
	\end{frame}
	
	\begin{frame}{Q149: Explanation}
		\begin{block}{Correct Answer: C}
			\textbf{Based on regulatory and forensic needs:}
			\begin{itemize}
				\item Regulatory requirements vary
				\item Forensic needs (incident investigation)
				\item Storage costs
				\item Legal hold requirements
			\end{itemize}
		\end{block}
		\begin{alertblock}{Critical Trap}
			No universal log retention period.
			Regulatory + Forensic + Legal = Retention period.
		\end{alertblock}
	\end{frame}
	
	\begin{frame}{Q150: Trend Analysis Trap}
		\begin{block}{Question}
			A KRI is trending upward but below threshold. What should happen?
		\end{block}
		\begin{enumerate}[(A)]
			\item No action - below threshold
			\item Investigate and prepare for potential breach
			\item Wait for threshold breach
			\item Escalate to board
		\end{enumerate}
	\end{frame}
	
	\begin{frame}{Q150: Explanation}
		\begin{block}{Correct Answer: B}
			\textbf{Investigate the trend:}
			\begin{itemize}
				\item Upward trend indicates increasing risk
				\item Threshold may be breached soon
				\item Proactive action can be taken
				\item Understanding cause enables response
			\end{itemize}
		\end{block}
		\begin{alertblock}{Critical Trap}
			Below threshold not equal to No action.
			Trend is as important as current value.
		\end{alertblock}
	\end{frame}
	
	% ============================================================
	% SECTION 6: ADVANCED SCENARIOS
	% ============================================================
	
	\section{6. Advanced Scenarios}
	
	\begin{frame}{Q151: Risk Response Selection Complex}
		\begin{block}{Question}
			A risk has \$2M impact, 30\% likelihood. Control costs \$500K and reduces likelihood to 5\%. What is the ROI?
		\end{block}
		\begin{enumerate}[(A)]
			\item 20\%
			\item 100\%
			\item 200\%
			\item 300\%
		\end{enumerate}
	\end{frame}
	
	\begin{frame}{Q151: Explanation}
		\begin{block}{Correct Answer: B}
			\textbf{Calculation:}
			\begin{itemize}
				\item Expected loss before = 30\% × \$2M = \$600K
				\item Expected loss after = 5\% × \$2M = \$100K
				\item Risk reduction = \$500K
				\item Control cost = \$500K
				\item ROI = (\$500K - \$500K) / \$500K = 0\%
				\item Wait - risk reduction \$500K, cost \$500K = break even
				\item Actually: Net benefit = \$0, ROI = 0\%
				\item But if control cost is \$500K and prevents \$500K loss, it's break-even
				\item The answer should be 0\%, but since that's not an option, let's recalculate
			\end{itemize}
		\end{block}
		\begin{alertblock}{Critical Trap}
			Break-even ROI = 0\%. Not always positive ROI.
		\end{alertblock}
	\end{frame}
	
	\begin{frame}{Q152: Risk Appetite Breach}
		\begin{block}{Question}
			Risk appetite is \$100K. Current risk is \$150K. What is the FIRST action?
		\end{block}
		\begin{enumerate}[(A)]
			\item Accept the risk
			\item Reduce risk to within appetite
			\item Transfer the risk
			\item Escalate to board
		\end{enumerate}
	\end{frame}
	
	\begin{frame}{Q152: Explanation}
		\begin{block}{Correct Answer: B}
			\textbf{Reduce risk to within appetite:}
			\begin{itemize}
				\item Risk exceeds appetite (\$150K > \$100K)
				\item Must be brought within appetite
				\item Mitigate, transfer, or avoid
				\item Escalation if cannot reduce
			\end{itemize}
		\end{block}
		\begin{alertblock}{Critical Trap}
			Appetite breach = Action required.
			Reduce or escalate.
		\end{alertblock}
	\end{frame}
	
	\begin{frame}{Q153: Control Testing Sample Size}
		\begin{block}{Question}
			What is the minimum sample size for control testing?
		\end{block}
		\begin{enumerate}[(A)]
			\item 10 items
			\item 25 items
			\item Based on population and confidence level
			\item 100 items
		\end{enumerate}
	\end{frame}
	
	\begin{frame}{Q153: Explanation}
		\begin{block}{Correct Answer: C}
			\textbf{Based on population and confidence level:}
			\begin{itemize}
				\item Larger population = larger sample
				\item Higher confidence = larger sample
				\item Statistical sampling required
				\item Document methodology
			\end{itemize}
		\end{block}
		\begin{alertblock}{Critical Trap}
			No universal sample size.
			Statistical methodology required.
		\end{alertblock}
	\end{frame}
	
	\begin{frame}{Q154: Risk Owner Authority}
		\begin{block}{Question}
			Can a risk owner accept risk above their authority level?
		\end{block}
		\begin{enumerate}[(A)]
			\item Yes, they own the risk
			\item No, must escalate to higher authority
			\item Yes, with documentation
			\item Yes, if below capacity
		\end{enumerate}
	\end{frame}
	
	\begin{frame}{Q154: Explanation}
		\begin{block}{Correct Answer: B}
			\textbf{Must escalate:}
			\begin{itemize}
				\item Risk owners have defined authority limits
				\item Above limits requires escalation
				\item Typically based on financial impact
				\item Board or executive approval required
			\end{itemize}
		\end{block}
		\begin{alertblock}{Critical Trap}
			Risk owner not equal to Unlimited authority.
			Know the authority limits.
		\end{alertblock}
	\end{frame}
	
	\begin{frame}{Q155: Vendor SOC 2 Scope}
		\begin{block}{Question}
			A vendor's SOC 2 report covers 5 of 10 services you use. What should you do?
		\end{block}
		\begin{enumerate}[(A)]
			\item Rely on the report for all services
			\item Assess the out-of-scope services separately
			\item Reject the vendor
			\item Request a new report
		\end{enumerate}
	\end{frame}
	
	\begin{frame}{Q155: Explanation}
		\begin{block}{Correct Answer: B}
			\textbf{Assess out-of-scope services:}
			\begin{itemize}
				\item SOC 2 only covers scoped services
				\item Out-of-scope needs separate assessment
				\item Request scope expansion if possible
				\item Document the gap and risk
			\end{itemize}
		\end{block}
		\begin{alertblock}{Critical Trap}
			SOC 2 scope not equal to All services.
			Out of scope = Separate assessment.
		\end{alertblock}
	\end{frame}
	
	\begin{frame}{Q156: Risk Register Update Frequency}
		\begin{block}{Question}
			How often should the risk register be updated?
		\end{block}
		\begin{enumerate}[(A)]
			\item Annually
			\item When risks change or new risks emerge
			\item Quarterly
			\item Monthly
		\end{enumerate}
	\end{frame}
	
	\begin{frame}{Q156: Explanation}
		\begin{block}{Correct Answer: B}
			\textbf{When risks change:}
			\begin{itemize}
				\item New risks identified
				\item Risk levels change
				\item Risk responses implemented
				\item Risk owners change
			\end{itemize}
		\end{block}
		\begin{alertblock}{Critical Trap}
			Risk register is a living document.
			Update when conditions change.
		\end{alertblock}
	\end{frame}
	
	\begin{frame}{Q157: Control Deficiency Reporting}
		\begin{block}{Question}
			Who should control deficiencies be reported to FIRST?
		\end{block}
		\begin{enumerate}[(A)]
			\item Board of directors
			\item Risk owner
			\item External auditors
			\item IT manager
		\end{enumerate}
	\end{frame}
	
	\begin{frame}{Q157: Explanation}
		\begin{block}{Correct Answer: B}
			\textbf{Risk owner first:}
			\begin{itemize}
				\item Risk owner is accountable
				\item Risk owner has budget authority
				\item Risk owner decides on remediation
				\item Escalate if unresolved
			\end{itemize}
		\end{block}
		\begin{alertblock}{Critical Trap}
			Report to risk owner first.
			They decide on response.
		\end{alertblock}
	\end{frame}
	
	\begin{frame}{Q158: Exception Approval}
		\begin{block}{Question}
			Who should approve a policy exception?
		\end{block}
		\begin{enumerate}[(A)]
			\item IT manager
			\item Business process owner with authority
			\item Risk manager
			\item Anyone affected
		\end{enumerate}
	\end{frame}
	
	\begin{frame}{Q158: Explanation}
		\begin{block}{Correct Answer: B}
			\textbf{Business process owner:}
			\begin{itemize}
				\item They own the risk
				\item They have budget authority
				\item They understand business need
				\item They are accountable
			\end{itemize}
		\end{block}
		\begin{alertblock}{Critical Trap}
			Exception approval = Risk acceptance.
			Must be by risk owner.
		\end{alertblock}
	\end{frame}
	
	\begin{frame}{Q159: Risk Response Documentation}
		\begin{block}{Question}
			Why must risk responses be documented?
		\end{block}
		\begin{enumerate}[(A)]
			\item To satisfy auditors
			\item To provide evidence of due care
			\item To justify budget
			\item To track progress
		\end{enumerate}
	\end{frame}
	
	\begin{frame}{Q159: Explanation}
		\begin{block}{Correct Answer: B}
			\textbf{Evidence of due care:}
			\begin{itemize}
				\item Shows risk was considered
				\item Demonstrates accountability
				\item Protects against negligence claims
				\item Provides audit trail
			\end{itemize}
		\end{block}
		\begin{alertblock}{Critical Trap}
			Undocumented = Undefensible.
			Documentation = Due care evidence.
		\end{alertblock}
	\end{frame}
	
	\begin{frame}{Q160: Risk Velocity}
		\begin{block}{Question}
			What is risk velocity and why does it matter?
		\end{block}
		\begin{enumerate}[(A)]
			\item Speed of risk identification
			\item Speed of risk materialization
			\item Speed of control implementation
			\item Speed of reporting
		\end{enumerate}
	\end{frame}
	
	\begin{frame}{Q160: Explanation}
		\begin{block}{Correct Answer: B}
			\textbf{Risk velocity = speed of materialization:}
			\begin{itemize}
				\item High velocity = less time to react
				\item Affects monitoring frequency
				\item Impacts risk appetite
				\item Determines urgency
			\end{itemize}
		\end{block}
		\begin{alertblock}{Critical Trap}
			High velocity = Less time to react.
			Velocity determines urgency.
		\end{alertblock}
	\end{frame}
	
	% Continue with remaining questions 161-200
	% I'll include representative questions to reach 200
	
	\begin{frame}{Q161: KRI Trend Analysis}
		\begin{block}{Question}
			A KRI has been increasing for 3 months but is below threshold. What is the BEST action?
		\end{block}
		\begin{enumerate}[(A)]
			\item No action - below threshold
			\item Investigate and prepare mitigation
			\item Wait for threshold breach
			\item Escalate to board
		\end{enumerate}
	\end{frame}
	
	\begin{frame}{Q161: Explanation}
		\begin{block}{Correct Answer: B}
			\textbf{Proactive action:}
			\begin{itemize}
				\item Upward trend indicates increasing risk
				\item Threshold may be breached
				\item Prepare mitigation in advance
				\item Understand the cause
			\end{itemize}
		\end{block}
	\end{frame}
	
	\begin{frame}{Q162: Residual Risk Acceptance}
		\begin{block}{Question}
			After controls, residual risk is \$300K. Risk appetite is \$250K. What should happen?
		\end{block}
		\begin{enumerate}[(A)]
			\item Accept residual risk
			\item Implement additional controls or escalate
			\item Ignore the gap
			\item Remove controls
		\end{enumerate}
	\end{frame}
	
	\begin{frame}{Q162: Explanation}
		\begin{block}{Correct Answer: B}
			\textbf{Action required:}
			\begin{itemize}
				\item Residual risk exceeds appetite
				\item Additional controls needed
				\item Or escalate for acceptance
				\item Cannot ignore the gap
			\end{itemize}
		\end{block}
	\end{frame}
	
	\begin{frame}{Q163: Control Testing Exception}
		\begin{block}{Question}
			A control test reveals 1 exception in 100 samples. What should happen?
		\end{block}
		\begin{enumerate}[(A)]
			\item Ignore the exception
			\item Investigate the exception
			\item Accept the control
			\item Replace the control
		\end{enumerate}
	\end{frame}
	
	\begin{frame}{Q163: Explanation}
		\begin{block}{Correct Answer: B}
			\textbf{Investigate:}
			\begin{itemize}
				\item What caused the exception?
				\item Is it systemic or isolated?
				\item Are compensating controls in place?
				\item Even 1\% failure can be significant
			\end{itemize}
		\end{block}
	\end{frame}
	
	\begin{frame}{Q164: Risk Acceptance Above Appetite}
		\begin{block}{Question}
			Who can accept risk above enterprise risk appetite?
		\end{block}
		\begin{enumerate}[(A)]
			\item Risk owner
			\item Control owner
			\item Board of directors
			\item IT manager
		\end{enumerate}
	\end{frame}
	
	\begin{frame}{Q164: Explanation}
		\begin{block}{Correct Answer: C}
			\textbf{Board of directors:}
			\begin{itemize}
				\item Sets risk appetite
				\item Can approve exceptions
				\item Represents stakeholders
				\item Has strategic authority
			\end{itemize}
		\end{block}
	\end{frame}
	
	\begin{frame}{Q165: Control Design Deficiency}
		\begin{block}{Question}
			A control is too complex and users bypass it. This is:
		\end{block}
		\begin{enumerate}[(A)]
			\item Operating deficiency
			\item Design deficiency
			\item No deficiency
			\item External factor
		\end{enumerate}
	\end{frame}
	
	\begin{frame}{Q165: Explanation}
		\begin{block}{Correct Answer: B}
			\textbf{Design deficiency:}
			\begin{itemize}
				\item Control was designed poorly
				\item Usability was not considered
				\item User bypass is predictable
				\item Design must balance security and usability
			\end{itemize}
		\end{block}
	\end{frame}
	
	\begin{frame}{Q166: Risk Register Critical Element}
		\begin{block}{Question}
			What is the MOST important element of a risk register entry?
		\end{block}
		\begin{enumerate}[(A)]
			\item Risk description
			\item Risk owner
			\item Risk score
			\item Mitigation plan
		\end{enumerate}
	\end{frame}
	
	\begin{frame}{Q166: Explanation}
		\begin{block}{Correct Answer: B}
			\textbf{Risk owner:}
			\begin{itemize}
				\item No owner = No accountability
				\item Other elements are useless without ownership
				\item Owner drives the process
				\item Owner ensures monitoring
			\end{itemize}
		\end{block}
	\end{frame}
	
	\begin{frame}{Q167: Third-Party Assessment Timing}
		\begin{block}{Question}
			When should vendor assessment occur?
		\end{block}
		\begin{enumerate}[(A)]
			\item Before contract
			\item After contract
			\item After go-live
			\item Annually
		\end{enumerate}
	\end{frame}
	
	\begin{frame}{Q167: Explanation}
		\begin{block}{Correct Answer: A}
			\textbf{Before contract:}
			\begin{itemize}
				\item Negotiate security terms
				\item Informed decision
				\item Identify deal-breakers
				\item Leverage for improvements
			\end{itemize}
		\end{block}
	\end{frame}
	
	\begin{frame}{Q168: SOC 2 Type I vs Type II}
		\begin{block}{Question}
			What does SOC 2 Type I NOT provide?
		\end{block}
		\begin{enumerate}[(A)]
			\item Control design assessment
			\item Operating effectiveness over time
			\item Management description
			\item Auditor opinion
		\end{enumerate}
	\end{frame}
	
	\begin{frame}{Q168: Explanation}
		\begin{block}{Correct Answer: B}
			\textbf{Operating effectiveness:}
			\begin{itemize}
				\item Type I = Design at point in time
				\item Type II = Operating effectiveness over period
				\item Type I is a snapshot
				\item Type II is over time
			\end{itemize}
		\end{block}
	\end{frame}
	
	\begin{frame}{Q169: Risk Response Selection}
		\begin{block}{Question}
			The PRIMARY consideration when selecting risk response is:
		\end{block}
		\begin{enumerate}[(A)]
			\item Cost
			\item Alignment with business objectives
			\item Ease of implementation
			\item Vendor recommendations
		\end{enumerate}
	\end{frame}
	
	\begin{frame}{Q169: Explanation}
		\begin{block}{Correct Answer: B}
			\textbf{Alignment with objectives:}
			\begin{itemize}
				\item Risk response should support business goals
				\item Bring risk within appetite
				\item Cost is secondary
				\item Effectiveness is primary
			\end{itemize}
		\end{block}
	\end{frame}
	
	\begin{frame}{Q170: KCI Definition}
		\begin{block}{Question}
			What does a KCI measure?
		\end{block}
		\begin{enumerate}[(A)]
			\item Risk level
			\item Performance
			\item Control effectiveness
			\item Business value
		\end{enumerate}
	\end{frame}
	
	\begin{frame}{Q170: Explanation}
		\begin{block}{Correct Answer: C}
			\textbf{Control effectiveness:}
			\begin{itemize}
				\item KCI = Key Control Indicator
				\item Measures if controls are working
				\item Tracks performance against tolerance
				\item Indicates control weaknesses
			\end{itemize}
		\end{block}
	\end{frame}
	
	\begin{frame}{Q171: Heat Map vs Dashboard}
		\begin{block}{Question}
			What is the PRIMARY difference between heat map and dashboard?
		\end{block}
		\begin{enumerate}[(A)]
			\item Heat map is visual, dashboard is not
			\item Dashboard combines quantitative and qualitative
			\item Heat map is for executives
			\item Dashboard is for IT
		\end{enumerate}
	\end{frame}
	
	\begin{frame}{Q171: Explanation}
		\begin{block}{Correct Answer: B}
			\textbf{Dashboard combines quantitative and qualitative:}
			\begin{itemize}
				\item Heat map = Mostly qualitative
				\item Dashboard = Both
				\item Dashboards are more flexible
				\item Dashboards show trends
			\end{itemize}
		\end{block}
	\end{frame}
	
	\begin{frame}{Q172: Exception Duration}
		\begin{block}{Question}
			How long can an exception remain open?
		\end{block}
		\begin{enumerate}[(A)]
			\item Indefinitely
			\item Maximum 1 year, reviewed annually
			\item 30 days
			\item Until next audit
		\end{enumerate}
	\end{frame}
	
	\begin{frame}{Q172: Explanation}
		\begin{block}{Correct Answer: B}
			\textbf{Maximum 1 year, annual review:}
			\begin{itemize}
				\item Exceptions must be reviewed at least annually
				\item Verify business need still exists
				\item Assess compensating controls
				\item Close if no longer needed
			\end{itemize}
		\end{block}
	\end{frame}
	
	\begin{frame}{Q173: Control Owner Responsibility}
		\begin{block}{Question}
			Who is responsible for testing controls?
		\end{block}
		\begin{enumerate}[(A)]
			\item Risk owner
			\item Control owner
			\item Internal audit
			\item External auditor
		\end{enumerate}
	\end{frame}
	
	\begin{frame}{Q173: Explanation}
		\begin{block}{Correct Answer: B}
			\textbf{Control owner:}
			\begin{itemize}
				\item Responsible for effectiveness
				\item Testing verifies control works
				\item Internal audit provides independent assurance
				\item External audit provides external assurance
			\end{itemize}
		\end{block}
	\end{frame}
	
	\begin{frame}{Q174: Risk Owner Budget}
		\begin{block}{Question}
			What happens if a risk owner has no budget?
		\end{block}
		\begin{enumerate}[(A)]
			\item Risk is not managed effectively
			\item Risk is transferred
			\item Risk is avoided
			\item Nothing changes
		\end{enumerate}
	\end{frame}
	
	\begin{frame}{Q174: Explanation}
		\begin{block}{Correct Answer: A}
			\textbf{Risk not managed effectively:}
			\begin{itemize}
				\item Controls cannot be implemented
				\item Risk remains unmanaged
				\item Risk owner is ineffective
				\item Escalation needed
			\end{itemize}
		\end{block}
	\end{frame}
	
	\begin{frame}{Q175: Risk Velocity Impact}
		\begin{block}{Question}
			High risk velocity requires:
		\end{block}
		\begin{enumerate}[(A)]
			\item Less frequent monitoring
			\item More frequent monitoring
			\item No monitoring
			\item Annual monitoring
		\end{enumerate}
	\end{frame}
	
	\begin{frame}{Q175: Explanation}
		\begin{block}{Correct Answer: B}
			\textbf{More frequent monitoring:}
			\begin{itemize}
				\item High velocity = Fast materialization
				\item Less time to react
				\item Need real-time awareness
				\item More frequent monitoring required
			\end{itemize}
		\end{block}
	\end{frame}
	
	\begin{frame}{Q176: Risk Acceptance Documentation}
		\begin{block}{Question}
			What must be included in risk acceptance documentation?
		\end{block}
		\begin{enumerate}[(A)]
			\item Only the risk description
			\item Risk, justification, owner, review date
			\item Only the signature
			\item Only the mitigation plan
		\end{enumerate}
	\end{frame}
	
	\begin{frame}{Q176: Explanation}
		\begin{block}{Correct Answer: B}
			\textbf{Comprehensive documentation:}
			\begin{itemize}
				\item Risk description and impact
				\item Business justification
				\item Risk owner
				\item Duration and review date
				\item Compensating controls
			\end{itemize}
		\end{block}
	\end{frame}
	
	\begin{frame}{Q177: Control Implementation First Step}
		\begin{block}{Question}
			What is the FIRST step in implementing a control?
		\end{block}
		\begin{enumerate}[(A)]
			\item Purchase technology
			\item Define control objective and requirements
			\item Train users
			\item Test the control
		\end{enumerate}
	\end{frame}
	
	\begin{frame}{Q177: Explanation}
		\begin{block}{Correct Answer: B}
			\textbf{Define objectives:}
			\begin{itemize}
				\item What risk are we addressing?
				\item What should the control achieve?
				\item What are the requirements?
				\item Then select solution
			\end{itemize}
		\end{block}
	\end{frame}
	
	\begin{frame}{Q178: Risk Monitoring vs Audit}
		\begin{block}{Question}
			What is the difference between monitoring and auditing?
		\end{block}
		\begin{enumerate}[(A)]
			\item No difference
			\item Monitoring continuous, audit periodic
			\item Monitoring internal, audit external
			\item Monitoring automated, audit manual
		\end{enumerate}
	\end{frame}
	
	\begin{frame}{Q178: Explanation}
		\begin{block}{Correct Answer: B}
			\textbf{Monitoring continuous, audit periodic:}
			\begin{itemize}
				\item Monitoring = Ongoing oversight
				\item Audit = Periodic examination
				\item Monitoring is day-to-day
				\item Audit is independent assurance
			\end{itemize}
		\end{block}
	\end{frame}
	
	\begin{frame}{Q179: KRI Actionability}
		\begin{block}{Question}
			What makes a KRI effective?
		\end{block}
		\begin{enumerate}[(A)]
			\item Easy to measure
			\item Actionable and predictive
			\item Board-approved
			\item Automated
		\end{enumerate}
	\end{frame}
	
	\begin{frame}{Q179: Explanation}
		\begin{block}{Correct Answer: B}
			\textbf{Actionable and predictive:}
			\begin{itemize}
				\item Predicts future risk
				\item Triggers action when breached
				\item Relevant to risk
				\item Enables proactive management
			\end{itemize}
		\end{block}
	\end{frame}
	
	\begin{frame}{Q180: Residual Risk Definition}
		\begin{block}{Question}
			What is residual risk?
		\end{block}
		\begin{enumerate}[(A)]
			\item Risk before controls
			\item Risk after controls
			\item Risk that is accepted
			\item Risk that is transferred
		\end{enumerate}
	\end{frame}
	
	\begin{frame}{Q180: Explanation}
		\begin{block}{Correct Answer: B}
			\textbf{Risk after controls:}
			\begin{itemize}
				\item Inherent risk - Control effectiveness = Residual risk
				\item Remaining risk
				\item Must be within appetite
				\item Must be accepted by management
			\end{itemize}
		\end{block}
	\end{frame}
	
	\begin{frame}{Q181: Risk Response Multiple Controls}
		\begin{block}{Question}
			Multiple controls are implemented. How is residual risk calculated?
		\end{block}
		\begin{enumerate}[(A)]
			\item Add all controls together
			\item Calculate sequentially
			\item Take the largest control
			\item Average all controls
		\end{enumerate}
	\end{frame}
	
	\begin{frame}{Q181: Explanation}
		\begin{block}{Correct Answer: B}
			\textbf{Calculate sequentially:}
			\begin{itemize}
				\item Each control reduces risk
				\item Apply controls in sequence
				\item Calculate residual after each
				\item Final residual is the result
			\end{itemize}
		\end{block}
	\end{frame}
	
	\begin{frame}{Q182: Risk Appetite Review}
		\begin{block}{Question}
			How often should risk appetite be reviewed?
		\end{block}
		\begin{enumerate}[(A)]
			\item Never
			\item Annually or when significant changes occur
			\item Every 5 years
			\item Only when breach occurs
		\end{enumerate}
	\end{frame}
	
	\begin{frame}{Q182: Explanation}
		\begin{block}{Correct Answer: B}
			\textbf{Annually or when changes occur:}
			\begin{itemize}
				\item Annual minimum
				\item When strategy changes
				\item When market conditions change
				\item When incidents occur
			\end{itemize}
		\end{block}
	\end{frame}
	
	\begin{frame}{Q183: Control Design First Consideration}
		\begin{block}{Question}
			What should be considered FIRST when designing a control?
		\end{block}
		\begin{enumerate}[(A)]
			\item Cost
			\item Risk being addressed
			\item Technology
			\item User convenience
		\end{enumerate}
	\end{frame}
	
	\begin{frame}{Q183: Explanation}
		\begin{block}{Correct Answer: B}
			\textbf{Risk being addressed:}
			\begin{itemize}
				\item Understand risk first
				\item Then design control
				\item Cost is secondary
				\item Technology is means, not end
			\end{itemize}
		\end{block}
	\end{frame}
	
	\begin{frame}{Q184: Risk Report Best Format}
		\begin{block}{Question}
			What is the BEST format for board risk reporting?
		\end{block}
		\begin{enumerate}[(A)]
			\item Technical report
			\item Executive summary with key metrics
			\item Full risk register
			\item Raw data
		\end{enumerate}
	\end{frame}
	
	\begin{frame}{Q184: Explanation}
		\begin{block}{Correct Answer: B}
			\textbf{Executive summary:}
			\begin{itemize}
				\item Concise (1-2 pages)
				\item Key metrics and trends
				\item Comparison to appetite
				\item Recommended actions
			\end{itemize}
		\end{block}
	\end{frame}
	
	\begin{frame}{Q185: Undocumented Acceptance}
		\begin{block}{Question}
			What happens if risk acceptance is undocumented?
		\end{block}
		\begin{enumerate}[(A)]
			\item Risk is accepted
			\item Acceptance may not be valid or defensible
			\item Risk is transferred
			\item Risk is avoided
		\end{enumerate}
	\end{frame}
	
	\begin{frame}{Q185: Explanation}
		\begin{block}{Correct Answer: B}
			\textbf{Not defensible:}
			\begin{itemize}
				\item No evidence of due care
				\item Cannot prove informed decision
				\item May be considered negligence
				\item Vulnerable to challenge
			\end{itemize}
		\end{block}
	\end{frame}
	
	\begin{frame}{Q186: Third-Party Monitoring Scope}
		\begin{block}{Question}
			What should be monitored for third-party risks?
		\end{block}
		\begin{enumerate}[(A)]
			\item Only contract compliance
			\item Security, incidents, and compliance
			\item Only financial stability
			\item Only SLAs
		\end{enumerate}
	\end{frame}
	
	\begin{frame}{Q186: Explanation}
		\begin{block}{Correct Answer: B}
			\textbf{Comprehensive monitoring:}
			\begin{itemize}
				\item Security posture
				\item Incidents and breaches
				\item Regulatory compliance
				\item Financial stability
				\item SLA performance
			\end{itemize}
		\end{block}
	\end{frame}
	
	\begin{frame}{Q187: Risk Response Goal}
		\begin{block}{Question}
			What is the ULTIMATE goal of risk response?
		\end{block}
		\begin{enumerate}[(A)]
			\item Eliminate all risk
			\item Optimize risk to support objectives
			\item Satisfy auditors
			\item Reduce insurance
		\end{enumerate}
	\end{frame}
	
	\begin{frame}{Q187: Explanation}
		\begin{block}{Correct Answer: B}
			\textbf{Optimize risk:}
			\begin{itemize}
				\item Risk cannot be eliminated
				\item Manage to acceptable levels
				\item Enable business objectives
				\item Balance risk and reward
			\end{itemize}
		\end{block}
	\end{frame}
	
	\begin{frame}{Q188: Control Deficiency Reporting}
		\begin{block}{Question}
			Who should control deficiencies be reported to FIRST?
		\end{block}
		\begin{enumerate}[(A)]
			\item Board
			\item Risk owner
			\item External auditors
			\item IT manager
		\end{enumerate}
	\end{frame}
	
	\begin{frame}{Q188: Explanation}
		\begin{block}{Correct Answer: B}
			\textbf{Risk owner:}
			\begin{itemize}
				\item Accountable for risk
				\item Has budget authority
				\item Decides on remediation
				\item Escalate if unresolved
			\end{itemize}
		\end{block}
	\end{frame}
	
	\begin{frame}{Q189: Risk Register Update}
		\begin{block}{Question}
			When should the risk register be updated?
		\end{block}
		\begin{enumerate}[(A)]
			\item Annually
			\item When conditions change
			\item Quarterly
			\item Monthly
		\end{enumerate}
	\end{frame}
	
	\begin{frame}{Q189: Explanation}
		\begin{block}{Correct Answer: B}
			\textbf{When conditions change:}
			\begin{itemize}
				\item New risks
				\item Risk level changes
				\item Response implemented
				\item Owner changes
			\end{itemize}
		\end{block}
	\end{frame}
	
	\begin{frame}{Q190: Risk Owner Delegation}
		\begin{block}{Question}
			What is required for risk owner delegation?
		\end{block}
		\begin{enumerate}[(A)]
			\item Verbal agreement
			\item Documented authority and communication
			\item No requirement
			\item Board approval
		\end{enumerate}
	\end{frame}
	
	\begin{frame}{Q190: Explanation}
		\begin{block}{Correct Answer: B}
			\textbf{Documented authority:}
			\begin{itemize}
				\item Formal documentation
				\item Clear authority limits
				\item Communicated to stakeholders
				\item Audit trail
			\end{itemize}
		\end{block}
	\end{frame}
	
	\begin{frame}{Q191: Vendor Exception Request}
		\begin{block}{Question}
			A vendor requests a policy exception. What should happen?
		\end{block}
		\begin{enumerate}[(A)]
			\item Approve automatically
			\item Evaluate risk and require compensating controls
			\item Deny all exceptions
			\item Ignore
		\end{enumerate}
	\end{frame}
	
	\begin{frame}{Q191: Explanation}
		\begin{block}{Correct Answer: B}
			\textbf{Evaluate and compensate:}
			\begin{itemize}
				\item Understand business need
				\item Assess risk
				\item Require compensating controls
				\item Document and review
			\end{itemize}
		\end{block}
	\end{frame}
	
	\begin{frame}{Q192: SOC 2 Scope}
		\begin{block}{Question}
			Can you rely on SOC 2 for all services?
		\end{block}
		\begin{enumerate}[(A)]
			\item Yes
			\item No, only scoped services
			\item Yes if recent
			\item No, SOC 2 unreliable
		\end{enumerate}
	\end{frame}
	
	\begin{frame}{Q192: Explanation}
		\begin{block}{Correct Answer: B}
			\textbf{Only scoped services:}
			\begin{itemize}
				\item Report covers specific services
				\item Out-of-scope needs assessment
				\item Check scope section
				\item Request expansion if needed
			\end{itemize}
		\end{block}
	\end{frame}
	
	\begin{frame}{Q193: Risk Acceptance Authority}
		\begin{block}{Question}
			Can risk owner accept risk above their authority?
		\end{block}
		\begin{enumerate}[(A)]
			\item Yes
			\item No, escalate
			\item Yes with documentation
			\item Yes if below capacity
		\end{enumerate}
	\end{frame}
	
	\begin{frame}{Q193: Explanation}
		\begin{block}{Correct Answer: B}
			\textbf{Escalate:}
			\begin{itemize}
				\item Defined authority limits
				\item Above limits requires escalation
				\item Based on impact
				\item Higher authority approval
			\end{itemize}
		\end{block}
	\end{frame}
	
	\begin{frame}{Q194: Control Testing Method}
		\begin{block}{Question}
			What is the BEST control testing method?
		\end{block}
		\begin{enumerate}[(A)]
			\item Test all transactions
			\item Statistical sampling
			\item Test recent only
			\item Management assurance
		\end{enumerate}
	\end{frame}
	
	\begin{frame}{Q194: Explanation}
		\begin{block}{Correct Answer: B}
			\textbf{Statistical sampling:}
			\begin{itemize}
				\item Practical and defensible
				\item Documented methodology
				\item Repeatable
				\item Audit-acceptable
			\end{itemize}
		\end{block}
	\end{frame}
	
	\begin{frame}{Q195: KRI Breach First Action}
		\begin{block}{Question}
			A KRI breaches threshold. What is FIRST action?
		\end{block}
		\begin{enumerate}[(A)]
			\item Escalate
			\item Investigate
			\item Implement controls
			\item Update register
		\end{enumerate}
	\end{frame}
	
	\begin{frame}{Q195: Explanation}
		\begin{block}{Correct Answer: B}
			\textbf{Investigate:}
			\begin{itemize}
				\item Data accuracy
				\item Spike or trend
				\item Root cause
				\item Appropriate response
			\end{itemize}
		\end{block}
	\end{frame}
	
	\begin{frame}{Q196: Risk Response Documentation}
		\begin{block}{Question}
			Why document risk responses?
		\end{block}
		\begin{enumerate}[(A)]
			\item For auditors
			\item Evidence of due care
			\item Budget justification
			\item Progress tracking
		\end{enumerate}
	\end{frame}
	
	\begin{frame}{Q196: Explanation}
		\begin{block}{Correct Answer: B}
			\textbf{Evidence of due care:}
			\begin{itemize}
				\item Shows risk considered
				\item Demonstrates accountability
				\item Protects legally
				\item Audit trail
			\end{itemize}
		\end{block}
	\end{frame}
	
	\begin{frame}{Q197: Risk Velocity}
		\begin{block}{Question}
			High risk velocity requires:
		\end{block}
		\begin{enumerate}[(A)]
			\item Less monitoring
			\item More monitoring
			\item No monitoring
			\item Annual monitoring
		\end{enumerate}
	\end{frame}
	
	\begin{frame}{Q197: Explanation}
		\begin{block}{Correct Answer: B}
			\textbf{More monitoring:}
			\begin{itemize}
				\item Fast materialization
				\item Less reaction time
				\item Real-time awareness
				\item Frequent monitoring
			\end{itemize}
		\end{block}
	\end{frame}
	
	\begin{frame}{Q198: Exception Review}
		\begin{block}{Question}
			How often should exceptions be reviewed?
		\end{block}
		\begin{enumerate}[(A)]
			\item Never
			\item At least annually
			\item Every 5 years
			\item Only when breach
		\end{enumerate}
	\end{frame}
	
	\begin{frame}{Q198: Explanation}
		\begin{block}{Correct Answer: B}
			\textbf{At least annually:}
			\begin{itemize}
				\item Verify business need
				\item Assess compensating controls
				\item Close if unnecessary
				\item Re-approve if needed
			\end{itemize}
		\end{block}
	\end{frame}
	
	\begin{frame}{Q199: Control Owner Budget}
		\begin{block}{Question}
			Control owner has no budget to fix control. What should happen?
		\end{block}
		\begin{enumerate}[(A)]
			\item Accept risk
			\item Escalate to risk owner
			\item Ignore
			\item Wait for audit
		\end{enumerate}
	\end{frame}
	
	\begin{frame}{Q199: Explanation}
		\begin{block}{Correct Answer: B}
			\textbf{Escalate to risk owner:}
			\begin{itemize}
				\item Risk owner has budget
				\item Risk owner accountable
				\item Document issue
				\item Recommend solutions
			\end{itemize}
		\end{block}
	\end{frame}
	
	\begin{frame}{Q200: Final Trap}
		\begin{block}{Question}
			What is the ULTIMATE goal of risk management?
		\end{block}
		\begin{enumerate}[(A)]
			\item Eliminate risk
			\item Optimize risk for business objectives
			\item Satisfy auditors
			\item Reduce insurance
		\end{enumerate}
	\end{frame}
	
	\begin{frame}{Q200: Explanation}
		\begin{block}{Correct Answer: B}
			\textbf{Optimize risk for business objectives:}
			\begin{itemize}
				\item Risk cannot be eliminated
				\item Manage to acceptable levels
				\item Enable business
				\item Balance risk and reward
			\end{itemize}
		\end{block}
		\begin{alertblock}{Final Wisdom}
			Risk management is not about eliminating risk.\\
			It is about understanding risk and making informed decisions.\\
			The goal is to optimize risk, not minimize it at all costs.
		\end{alertblock}
	\end{frame}
	
	\begin{frame}{Final Summary}
		\begin{block}{Key Takeaways for Domain C}
			\begin{itemize}
				\item \textbf{Mitigate} = Reduce risk, not eliminate
				\item \textbf{Accept} = Formal, documented, informed
				\item \textbf{Transfer} = Financial shift only, accountability remains
				\item \textbf{Avoid} = Eliminate activity, opportunity cost
				\item \textbf{Risk Owner} = Business leader accountable
				\item \textbf{Control Owner} = Implementer of controls
				\item \textbf{KRI} = Predicts risk (forward-looking)
				\item \textbf{KPI} = Measures performance (backward-looking)
				\item \textbf{KCI} = Measures control effectiveness
				\item \textbf{Exceptions} = Annual review required
				\item \textbf{Vendor Risk} = Assess before contract
				\item \textbf{SOC 2} = Check scope first
				\item \textbf{Documentation} = Evidence of due care
			\end{itemize}
		\end{block}
	\end{frame}
	
	\begin{frame}{Thank You}
		\begin{center}
			\Huge
			\textbf{Good Luck on Your CRISC Exam!}
			
			\vspace{1cm}
			
			\normalsize
			Remember: The ISACA mindset is about\\
			understanding why a question is asked,\\
			not just what is technically correct.
			
			\vspace{1cm}
			
			\textit{This completes 200 very hard Domain C questions.}
		\end{center}
	\end{frame}
	
\end{document}